\documentclass[10pt,a4paper]{article}
\usepackage{fontspec,amsmath,amssymb,graphicx}
\usepackage[margin=18mm]{geometry}
\setmainfont{Times New Roman}
\setlength{\parindent}{0pt}
\setlength{\parskip}{5pt}
\allowdisplaybreaks
\sloppy
\emergencystretch=2em
\begin{document}
\section*{BT\_4\_2\_3\_02 — Lời giải bổ sung}
\subsection*{Kiến thức và trình tự giải}
\textbf{Bước 1 — Viết toán tử và biên parabol.} Đặt $P=\partial_t-a(x,t)\partial_{xx}-b(x,t)\partial_x-c(x,t)$ với $a\ge0$, $c\le0$. Biên parabol gồm mặt $t=0$ và các biên không gian; mặt $t=T$ vẫn thuộc miền được ước lượng. Với phương trình nhiệt, $a=K>0$, $b=c=0$.

\textbf{Bước 2 — Loại cực đại dương bằng dấu đạo hàm.} Tại cực đại dương nội miền, $w_x=0$, $w_{xx}\le0$, $w_t=0$; nếu cực đại trên mặt thời gian trên thì đạo hàm một phía cho $w_t\ge0$. Suy $Pw\ge0$, mâu thuẫn $Pw<0$. Với $Pw\le0$, xét $w_\varepsilon=w-\varepsilon e^{\gamma t}$, $\gamma>0$: $Pw_\varepsilon=Pw-\varepsilon(\gamma-c)e^{\gamma t}<0$. Chứng minh $w_\varepsilon\le0$ rồi cho $\varepsilon\downarrow0$.

\textbf{Bước 3 — So sánh hai hàm bằng hiệu.} Để chứng minh $v\le u\le V$, kiểm tra $Pv\le Pu\le PV$ và thứ tự của dữ liệu đầu, biên. Áp dụng bước 2 cho $v-u$ và $u-V$. Nếu có nguồn dương, hằng số không nhất thiết là nghiệm trên, nên phải tính $PV-Pu$; không dùng cực đại của dữ liệu để chặn một bài có nguồn mà chưa kiểm tra.

\textbf{Bước 4 — Xử lý Neumann, Robin và tuần hoàn.} Neumann thuần nhất có thể phản xạ chẵn qua đầu mút; điểm biên trở thành điểm trong. Robin viết theo pháp tuyến ngoài $\partial_\nu u+\sigma u=0$, $\sigma\ge0$. Tại cực đại dương không hằng, bổ đề Hopf cho $\partial_\nu u>0$, trái với $\partial_\nu u=-\sigma u\le0$. Với Robin phải thêm mức $0$ vào cận hằng nếu dữ liệu mang một dấu. Biên tuần hoàn ghép hai đầu thành đường tròn; không còn biên không gian.

\textbf{Bước 5 — Miền vô hạn và nguồn bị chặn.} Trên toàn trục/nửa trục, cần nghiệm bị chặn trên mỗi dải thời gian (hoặc điều kiện tăng trưởng đủ dùng). Dùng rào $\varepsilon((x-x_0)^2+2Kt)$ trên đoạn lớn, áp dụng cực đại trên miền hữu hạn rồi cho $\varepsilon\downarrow0$; cách này loại ảnh hưởng từ vô cực. Nếu $|Pu|\le N$ và $|u|\le M$ trên biên parabol, đặt $V=M+Nt$. Vì $PV=N-cV\ge N$, suy $|u|\le M+Nt\le M+NT$. $M$ ở đây là cận của $|u|$, không phải cực đại có dấu.
\subsection*{Lời giải áp dụng cho bài}
\textbf{Bước 1. Phát biểu trên nửa trục.} Với nghiệm bị chặn trên $[0,\infty)\times[0,T]$ và $f=u(\cdot,0)$:
\begin{align*}
(a)&:\quad \min(\inf_{x\ge0}f,\inf_{0\le s\le T}\mu(s))
\le u\le\max(\sup_{x\ge0}f,\sup_s\mu(s));\\
(b)&:\quad \inf_{x\ge0}f\le u\le\sup_{x\ge0}f.
\end{align*}
\textbf{Bước 2. Dirichlet và vô cực.} Gọi cận trên (a) là $C$, đặt $w=u-C\le0$ trên mặt đầu và tại $x=0$. Xét $w_\varepsilon=w-\varepsilon(x^2+2Kt)$: vẫn giải PDE, không dương trên hai phần biên này. Vì $w$ bị chặn, chọn $R$ lớn làm $w_\varepsilon(R,t)\le0$ cho $0\le t\le T$. Cực đại trên $(0,R)$ cho $w_\varepsilon\le0$. Với $(x,t)$ cố định, cho $\varepsilon\downarrow0$ được $w\le0$. Áp dụng cho $-u$ được cận dưới.
\textbf{Bước 3. Neumann.} Mở rộng chẵn $\widetilde u(x,t)=u(|x|,t)$. Biên $u_x(0,t)=0$ làm nghiệm mở rộng trơn qua $0$ với $t>0$, bị chặn trên toàn trục. Cận toàn trục bằng rào bậc hai cho $\inf f\le\widetilde u\le\sup f$, hạn chế lại $x\ge0$ là kết quả (b).
Miền in $\mathbb R$ ở dòng PDE không khớp phần mở đầu nửa trục; chứng minh dùng miền $x>0$ và biên được nêu trong bài.\end{document}